\capitulosinnumero{Resumen}

Este documento se corresponde con el Trabajo de Fin de Grado del Grado en Ingeniería Multimedia de Leigh Garrett. El presente trabajo, titulado «Desarrollo de un videojuego narrativo en 2D con mecánicas adaptativas sobre la perspectiva de género» y realizado en el curso académico 2026-2027, ha sido supervisado y tutorizado por Carlos J. Villagrá Arnedo y Patricia María Compañ Rosique.
\\\\
En los últimos años, el videojuego se ha consolidado como un medio artístico y expresivo capaz de abordar temáticas profundamente humanas e introspectivas. El presente proyecto consiste en la conceptualización, diseño y desarrollo de un videojuego narrativo en 2D centrado en el agotamiento psicológico derivado de una crisis de identidad de género a causa del rechazo social e internalizado. A partir de la perspectiva de género, la obra emplea la metáfora de un escenario apocalíptico y la exploración de una serie de mundos abstractos para representar el esfuerzo constante por encajar en las expectativas normativas de la sociedad.
\\\\
La principal aportación técnica y de diseño del proyecto radica en la implementación de un sistema de jugabilidad dinámica (gameplay adaptativo) en segundo plano. Este módulo analiza continuamente el comportamiento y las acciones del jugador durante la partida (identificando tendencias agresivas, pasivas o evasivas) con el fin de modificar dinámicamente las mecánicas, las habilidades del protagonista e incluso la propia interfaz gráfica de usuario (UI), reflejando así el estado psicológico interno del personaje. Con el propósito de evitar la frustración y garantizar la agencia del usuario, dicho sistema solicita el consentimiento explícito del jugador antes de aplicar las transformaciones en la experiencia de juego. Asimismo, para tratar estas cuestiones con el rigor y la sensibilidad necesarios, el proyecto contempla el asesoramiento de especialistas en el área de la psicología.
\\\\
El objetivo final del juego es sobrevivir a este viaje introspectivo hasta alcanzar un desenlace abierto, cuya naturaleza e interpretación dependerán directamente del vínculo construido entre el jugador y el protagonista a través de su estilo de juego. Esta memoria abarca desde el estudio del estado del arte y la fundamentación teórica sobre narrativa digital y mecánicas adaptativas, hasta la elaboración del Documento de Diseño de Juego (GDD), la dirección artística, la arquitectura técnica en el motor Unity y el desarrollo de un prototipo funcional y jugable.
\\\\
Palabras clave: Videojuego 2D, Narrativa Interactiva, Mecánicas Adaptativas, Perspectiva de Género, Salud Mental, Identidad de Género, Unity.